% Généré pour inclusion dans rapport-tp-ldap-guide-operateur.tex - ne pas éditer à la main sans resynchroniser le dépôt.
\section{\texttt{docker-compose.yml}}\label{ann:compose}
{\footnotesize
\begin{verbatim}
name: tp-ldap

services:
  ldap:
    build:
      context: .
      dockerfile: docker/Dockerfile
    container_name: ldap
    environment:
      LDAP_ORGANISATION: Example Org
      LDAP_DOMAIN: example.org
      LDAP_BASE_DN: dc=example,dc=org
      LDAP_ADMIN_PASSWORD: admin
      LDAP_TLS: "false"
    ports:
      - "389:389"
    volumes:
      - ldap_data:/var/lib/ldap
      - ldap_config:/etc/ldap/slapd.d
    restart: "no"
    stop_grace_period: 15s
    healthcheck:
      # Fichier écrit après tous les scripts init.d ; puis DIT + syncprov (réplication).
      test:
        [
          "CMD-SHELL",
          "test -f /container/.ldap-init-complete && ldapsearch -x -H ldap://127.0.0.1:389 -D cn=admin,dc=example,dc=org -w \"$$LDAP_ADMIN_PASSWORD\" -b dc=example,dc=org -s base dn >/dev/null 2>&1 && ldapsearch -H ldapi:/// -Y EXTERNAL -b cn=config -s sub -LLL '(objectClass=olcSyncProvConfig)' dn 2>/dev/null | grep -q '^dn:'",
        ]
      interval: 5s
      timeout: 5s
      retries: 12
      start_period: 30s

  ldap-replica:
    build:
      context: .
      dockerfile: docker/Dockerfile
    container_name: ldap-replica
    environment:
      LDAP_ORGANISATION: Example Org
      LDAP_DOMAIN: example.org
      LDAP_BASE_DN: dc=example,dc=org
      LDAP_ADMIN_PASSWORD: admin
      LDAP_TLS: "false"
      LDAP_SERVICE_ROLE: consumer
      LDAP_REPLICATION_PROVIDER_URI: ldap://ldap:389
    ports:
      - "1389:389"
    volumes:
      - ldap_replica_data:/var/lib/ldap
      - ldap_replica_config:/etc/ldap/slapd.d
    depends_on:
      ldap:
        condition: service_healthy
    healthcheck:
      test:
        [
          "CMD-SHELL",
          "test -f /container/.ldap-init-complete && ldapsearch -x -H ldap://127.0.0.1:389 -b '' -s base namingContexts >/dev/null 2>&1",
        ]
      interval: 5s
      timeout: 5s
      retries: 15
      start_period: 60s
    restart: "no"
    stop_grace_period: 15s

  # Second annuaire (entité distincte) pour le méta-annuaire - même image, autre suffixe.
  ldap-acme:
    build:
      context: .
      dockerfile: docker/Dockerfile
    container_name: ldap-acme
    environment:
      LDAP_ORGANISATION: Acme Corp
      LDAP_DOMAIN: acme.com
      LDAP_BASE_DN: dc=acme,dc=com
      LDAP_ADMIN_PASSWORD: admin
      LDAP_TLS: "false"
    ports:
      - "2389:389"
    volumes:
      - ldap_acme_data:/var/lib/ldap
      - ldap_acme_config:/etc/ldap/slapd.d
    depends_on:
      ldap:
        condition: service_healthy
    healthcheck:
      test:
        [
          "CMD-SHELL",
          "test -f /container/.ldap-init-complete && ldapsearch -x -H ldap://127.0.0.1:389 -D cn=admin,dc=acme,dc=com -w \"$$LDAP_ADMIN_PASSWORD\" -b dc=acme,dc=com -s base dn >/dev/null 2>&1 && ldapsearch -H ldapi:/// -Y EXTERNAL -b cn=config -s sub -LLL '(objectClass=olcSyncProvConfig)' dn 2>/dev/null | grep -q '^dn:'",
        ]
      interval: 5s
      timeout: 5s
      retries: 12
      start_period: 30s
    restart: "no"
    stop_grace_period: 15s

  # Méta-annuaire (back-meta) : agrège ldap (example.org) et ldap-acme (acme.com).
  ldap-meta:
    build:
      context: .
      dockerfile: docker/Dockerfile
    container_name: ldap-meta
    environment:
      LDAP_SERVICE_ROLE: meta
      LDAP_ADMIN_PASSWORD: admin
      LDAP_TLS: "false"
      LDAP_META_SUFFIX: o=federation
      LDAP_META_EXAMPLE_URI: ldap://ldap:389
      LDAP_META_EXAMPLE_SUFFIX: dc=example,dc=org
      LDAP_META_EXAMPLE_VIRTUAL: ou=example-org,o=federation
      LDAP_META_ACME_URI: ldap://ldap-acme:389
      LDAP_META_ACME_SUFFIX: dc=acme,dc=com
      LDAP_META_ACME_VIRTUAL: ou=acme-corp,o=federation
    ports:
      - "3389:389"
    volumes:
      - ldap_meta_data:/var/lib/ldap
      - ldap_meta_config:/etc/ldap/slapd.d
    depends_on:
      ldap:
        condition: service_healthy
      ldap-acme:
        condition: service_healthy
    healthcheck:
      test:
        [
          "CMD-SHELL",
          "test -f /container/.ldap-init-complete && ldapsearch -x -H ldap://127.0.0.1:389 -D cn=admin,o=federation -w \"$$LDAP_ADMIN_PASSWORD\" -b uid=thomas,ou=people,ou=example-org,o=federation -s base dn >/dev/null 2>&1 && ldapsearch -x -H ldap://127.0.0.1:389 -D cn=admin,o=federation -w \"$$LDAP_ADMIN_PASSWORD\" -b uid=thomas,ou=people,ou=acme-corp,o=federation -s base dn >/dev/null 2>&1",
        ]
      interval: 5s
      timeout: 5s
      retries: 15
      start_period: 45s
    restart: "no"
    stop_grace_period: 15s

  keycloak:
    image: quay.io/keycloak/keycloak:26.2
    container_name: keycloak
    command: ["start-dev", "--http-port=8080", "--hostname-strict=false"]
    environment:
      KC_BOOTSTRAP_ADMIN_USERNAME: admin
      KC_BOOTSTRAP_ADMIN_PASSWORD: admin
      KC_HEALTH_ENABLED: "true"
      KC_HTTP_MANAGEMENT_PORT: "9000"
    ports:
      # 8090 / 9001 : éviter le conflit avec un service déjà sur 8080 (ex. autre stack).
      - "8090:8080"
      - "9001:9000"
    depends_on:
      ldap:
        condition: service_healthy
    restart: "no"

volumes:
  ldap_data: {}
  ldap_config: {}
  ldap_replica_data: {}
  ldap_replica_config: {}
  ldap_acme_data: {}
  ldap_acme_config: {}
  ldap_meta_data: {}
  ldap_meta_config: {}
\end{verbatim}
}

\section{\texttt{docker/Dockerfile}}\label{ann:dockerfile}
{\footnotesize
\begin{verbatim}
FROM debian:12-slim

ENV DEBIAN_FRONTEND=noninteractive \
    LDAP_ORGANISATION="Example Org" \
    LDAP_DOMAIN="example.org" \
    LDAP_BASE_DN="dc=example,dc=org" \
    LDAP_ADMIN_PASSWORD="admin" \
    LDAP_TLS="false"

# slapd + clients ; NSS/PAM + nslcd/nscd pour l’objectif 4 (getent).
RUN apt-get update \
    && apt-get install -y --no-install-recommends \
       slapd ldap-utils tini bash coreutils procps ca-certificates sysvinit-utils \
       libpam-ldap libnss-ldap nslcd nscd \
    && rm -rf /var/lib/apt/lists/*

RUN mkdir -p /container/init.d /var/lib/ldap /etc/ldap/slapd.d \
    && chown -R openldap:openldap /var/lib/ldap /etc/ldap/slapd.d

COPY --chown=openldap:openldap scripts/init_ldap.sh /container/init.d/init_ldap.sh
COPY --chown=openldap:openldap scripts/init_ldap_linux_integration.sh /container/init.d/init_ldap_linux_integration.sh
COPY --chown=openldap:openldap scripts/init_replication_provider.sh /container/init.d/init_replication_provider.sh
COPY --chown=openldap:openldap scripts/init_replication_consumer.sh /container/init.d/init_replication_consumer.sh
COPY --chown=root:root scripts/init_meta_annuaire.sh /container/init_meta_annuaire.sh
COPY docker/entrypoint.sh /entrypoint.sh
RUN chmod +x /entrypoint.sh /container/init_meta_annuaire.sh /container/init.d/init_ldap.sh /container/init.d/init_ldap_linux_integration.sh \
    /container/init.d/init_replication_provider.sh /container/init.d/init_replication_consumer.sh

EXPOSE 389

VOLUME ["/var/lib/ldap", "/etc/ldap/slapd.d"]

ENTRYPOINT ["/usr/bin/tini", "--", "/entrypoint.sh"]
\end{verbatim}
}

\section{\texttt{docker/entrypoint.sh}}\label{ann:entrypoint}
{\footnotesize
\begin{verbatim}
#!/usr/bin/env bash
set -euo pipefail

# Script d'entrée pour le conteneur LDAP
# Ce script démarre slapd et exécute les scripts d'initialisation

echo "[entrypoint] Démarrage du conteneur LDAP..."

# Les volumes Docker montés sur /var/lib/ldap et /etc/ldap/slapd.d sont vides
# et souvent propriété de root : slapd (-u openldap) ne peut pas y accéder.
chown -R openldap:openldap /var/lib/ldap /etc/ldap/slapd.d

# Initialisation de la configuration slapd si nécessaire
# Cette étape crée la structure de base cn=config
if [ ! -d "/etc/ldap/slapd.d" ] || [ -z "$(ls -A /etc/ldap/slapd.d || true)" ]; then
  echo "[entrypoint] Initialisation de la configuration slapd..."
  slaptest -F /etc/ldap/slapd.d -f /dev/null >/dev/null 2>&1 || true
  chown -R openldap:openldap /etc/ldap/slapd.d
fi

# Vérifier si slapd est déjà en cours d'exécution
if pgrep slapd >/dev/null 2>&1; then
    echo "[entrypoint] slapd déjà en cours d'exécution"
else
    # Démarrage du serveur slapd en arrière-plan.
    # slapd se daemonise (fork) : le PID du shell ($!) meurt tout de suite, on vérifie la vie du
    # service avec pgrep slapd, pas avec kill -0 $!.
    echo "[entrypoint] Lancement de slapd..."
    slapd -h "ldap:/// ldapi:///" -u openldap -g openldap -F /etc/ldap/slapd.d &
fi

# Attendre que slapd écoute (processus slapd + socket ldapi + LDAP TCP).
echo "[entrypoint] Attente du socket ldapi et du port LDAP 389..."
ready=0
for i in {1..60}; do
  if ! pgrep slapd >/dev/null 2>&1; then
    echo "[entrypoint] ERREUR: aucun processus slapd. Vérifiez les droits sur /etc/ldap/slapd.d et /var/lib/ldap (propriétaire openldap)."
    exit 1
  fi
  if { [ -S /var/run/slapd/ldapi ] || [ -S /run/slapd/ldapi ]; } && bash -c ":>/dev/tcp/127.0.0.1/389" 2>/dev/null; then
    echo "[entrypoint] slapd prêt (ldapi + port 389)"
    ready=1
    break
  fi
  sleep 1
done
if [ "$ready" != 1 ]; then
  echo "[entrypoint] ERREUR: slapd ne semble pas joignable après 60 s."
  exit 1
fi

# Exécution des scripts d'initialisation (ordre fixe : pas de préfixe numérique sur les noms)
if [ -d "/container/init.d" ]; then
  echo "[entrypoint] Exécution des scripts d'initialisation..."
  for name in init_ldap.sh init_ldap_linux_integration.sh init_replication_provider.sh init_replication_consumer.sh; do
    script="/container/init.d/$name"
    [ -e "$script" ] || continue
    echo "[entrypoint] Exécution: $script"
    if bash "$script"; then
      echo "[entrypoint] Terminé OK: $script"
    else
      rc=$?
      echo "[entrypoint] AVERTISSEMENT: $script a retourné le code $rc (poursuite du démarrage)"
    fi
  done
  echo "[entrypoint] Scripts d'initialisation terminés"
  touch /container/.ldap-init-complete
fi

# Arrêt propre : docker stop envoie SIGTERM à PID 1
# Éviter « sleep infinity » en boucle : l'arrêt peut rester bloqué très longtemps en « Stopping »
_shutdown() {
  echo "[entrypoint] Arrêt demandé (SIGTERM/SIGINT)..."
  pkill -u openldap -TERM slapd 2>/dev/null || pkill -TERM slapd 2>/dev/null || true
  local i
  for i in $(seq 1 40); do
    pgrep slapd >/dev/null 2>&1 || break
    sleep 0.25
  done
  pkill -u openldap -KILL slapd 2>/dev/null || pkill -KILL slapd 2>/dev/null || true
  exit 0
}
trap _shutdown TERM INT

echo "[entrypoint] Conteneur prêt (docker stop pour arrêter)..."
tail -f /dev/null &
wait $!
\end{verbatim}
}

\section{Script \texttt{scripts/init\_ldap.sh}}\label{ann:init-ldap}
{\footnotesize
\begin{verbatim}
#!/usr/bin/env bash
# Script d'initialisation LDAP
# Ce script configure la base mdb, crée le DIT et configure les ACL

# Désactiver l'arrêt sur erreur pour permettre la continuation
set -uo pipefail

# Variables d'environnement attendues:
# - LDAP_BASE_DN
# - LDAP_ORGANISATION
# - LDAP_DOMAIN
# - LDAP_ADMIN_PASSWORD (utilisée pour le rootDN de la base)
# - LDAP_SERVICE_ROLE : provider (défaut) = DIT complet ; consumer = base mdb + ACL, données via syncrepl

BASE_DN=${LDAP_BASE_DN:-"dc=polytech,dc=fr"}
ROLE=${LDAP_SERVICE_ROLE:-provider}
ORG=${LDAP_ORGANISATION:-"Polytech"}
DOMAIN=${LDAP_DOMAIN:-"polytech.fr"}
ADMIN_PASS=${LDAP_ADMIN_PASSWORD:-"admin123"}

export LDAPTLS_REQCERT=never # pour ne pas avoir à fournir de certificat lors de la connexion

# Attendre que slapd réponde (ldapi)
for i in {1..30}; do
  if ldapwhoami -H ldapi:/// -Y EXTERNAL >/dev/null 2>&1; then
    break
  fi
  sleep 1
done

if [ "$ROLE" = "meta" ]; then
  echo "[init_ldap] Rôle meta : méta-annuaire (pas de DIT mdb local)."
  bash /container/init_meta_annuaire.sh
  exit $?
fi

# Supprimer la base mdb par défaut et créer une nouvelle
echo "[init_ldap] Suppression de la base mdb par défaut..."
ldapmodify -H ldapi:/// -Y EXTERNAL <<EOF
dn: olcDatabase={1}mdb,cn=config
changetype: delete
EOF
echo "[init_ldap] Base mdb par défaut supprimée"

# Les fichiers sous /var/lib/ldap gardent sinon l’ancien suffixe (ex. dc=nodomain) alors que
# cn=config pointe déjà vers LDAP_BASE_DN → syncrepl / slapcat incohérents.
echo "[init_ldap] Purge des fichiers mdb et redémarrage de slapd..."
pkill -u openldap -TERM slapd 2>/dev/null || pkill -TERM slapd 2>/dev/null || true
for _ in $(seq 1 40); do
  pgrep slapd >/dev/null 2>&1 || break
  sleep 0.25
done
pkill -u openldap -KILL slapd 2>/dev/null || pkill -KILL slapd 2>/dev/null || true
find /var/lib/ldap -mindepth 1 -delete 2>/dev/null || true
chown openldap:openldap /var/lib/ldap
slapd -h "ldap:/// ldapi:///" -u openldap -g openldap -F /etc/ldap/slapd.d &
for _ in $(seq 1 60); do
  if ldapwhoami -H ldapi:/// -Y EXTERNAL >/dev/null 2>&1; then
    break
  fi
  sleep 1
done

# Créer une nouvelle base mdb avec la configuration correcte
echo "[init_ldap] Création de la base mdb pour suffix=$BASE_DN..."
HASH=$(slappasswd -s "$ADMIN_PASS")
cat > /tmp/new-db.ldif <<EOF
dn: olcDatabase=mdb,cn=config
objectClass: olcDatabaseConfig
objectClass: olcMdbConfig
olcDatabase: mdb
olcSuffix: $BASE_DN
olcRootDN: cn=admin,$BASE_DN
olcRootPW: $HASH
olcDbDirectory: /var/lib/ldap
olcDbIndex: objectClass eq
olcDbIndex: cn,sn,uid eq
olcDbMaxSize: 1073741824
EOF
ldapadd -H ldapi:/// -Y EXTERNAL -f /tmp/new-db.ldif
echo "[init_ldap] Base mdb créée avec suffix=$BASE_DN"

if [ "$ROLE" = "consumer" ]; then
  echo "[init_ldap] Rôle consumer : pas de création locale du DIT (réplication depuis le fournisseur)."
else
# DIT: base + ou=people + ou=groups + groupe admin_ldap + groupe developers + utilisateurs thomas et john
cat > /tmp/dit.ldif <<EOF
version: 1

dn: $BASE_DN
objectClass: top
objectClass: dcObject
objectClass: organization
o: $ORG
dc: ${DOMAIN%%.*}

dn: ou=people,$BASE_DN
objectClass: top
objectClass: organizationalUnit
ou: people

dn: ou=groups,$BASE_DN
objectClass: top
objectClass: organizationalUnit
ou: groups

dn: cn=admin_ldap,ou=groups,$BASE_DN
objectClass: top
objectClass: groupOfNames
cn: admin_ldap
member: cn=dummy,$BASE_DN

dn: cn=developers,ou=groups,$BASE_DN
objectClass: top
objectClass: groupOfNames
cn: developers
member: cn=dummy,$BASE_DN

dn: cn=admin_keycloak,ou=groups,$BASE_DN
objectClass: top
objectClass: groupOfNames
cn: admin_keycloak
member: cn=dummy,$BASE_DN

dn: uid=thomas,ou=people,$BASE_DN
objectClass: top
objectClass: inetOrgPerson
cn: Thomas
sn: Thomas
uid: thomas
mail: thomas@$DOMAIN
userPassword: $(slappasswd -s thomas123)

dn: uid=john,ou=people,$BASE_DN
objectClass: top
objectClass: inetOrgPerson
cn: John
sn: John
uid: john
mail: john@$DOMAIN
userPassword: $(slappasswd -s john123)
EOF

# Ajouter la structure si absente (bind simple sur rootDN)
echo "[init_ldap] Vérification de l'existence du DIT..."
if ! ldapsearch -x -H ldap:/// -D "cn=admin,$BASE_DN" -w "$ADMIN_PASS" -b "$BASE_DN" -s base dn >/dev/null 2>&1; then
  echo "[init_ldap] Création DIT de base..."
  ldapadd -x -H ldap:/// -D "cn=admin,$BASE_DN" -w "$ADMIN_PASS" -f /tmp/dit.ldif || echo "[init_ldap] Erreur lors de la création du DIT"
  echo "[init_ldap] DIT créé"
else
  echo "[init_ldap] DIT déjà présent"
fi
echo "[init_ldap] Étape DIT terminée"

# Ajouter thomas au groupe admin_ldap et john au groupe developers
echo "[init_ldap] Ajout de thomas au groupe admin_ldap..."
cat > /tmp/add-thomas-to-admin.ldif <<EOF
dn: cn=admin_ldap,ou=groups,$BASE_DN
changetype: modify
delete: member
member: cn=dummy,$BASE_DN
-
add: member
member: uid=thomas,ou=people,$BASE_DN
EOF
ldapmodify -x -H ldap:/// -D "cn=admin,$BASE_DN" -w "$ADMIN_PASS" -f /tmp/add-thomas-to-admin.ldif || echo "[init_ldap] Erreur lors de l'ajout de thomas au groupe admin_ldap"
echo "[init_ldap] Thomas ajouté au groupe admin_ldap"

echo "[init_ldap] Ajout de john au groupe developers..."
cat > /tmp/add-john-to-developers.ldif <<EOF
dn: cn=developers,ou=groups,$BASE_DN
changetype: modify
delete: member
member: cn=dummy,$BASE_DN
-
add: member
member: uid=john,ou=people,$BASE_DN
EOF
ldapmodify -x -H ldap:/// -D "cn=admin,$BASE_DN" -w "$ADMIN_PASS" -f /tmp/add-john-to-developers.ldif || echo "[init_ldap] Erreur lors de l'ajout de john au groupe developers"
echo "[init_ldap] John ajouté au groupe developers"

echo "[init_ldap] Ajout de thomas au groupe admin_keycloak..."
cat > /tmp/add-thomas-to-admin-keycloak.ldif <<EOF
dn: cn=admin_keycloak,ou=groups,$BASE_DN
changetype: modify
delete: member
member: cn=dummy,$BASE_DN
-
add: member
member: uid=thomas,ou=people,$BASE_DN
EOF
ldapmodify -x -H ldap:/// -D "cn=admin,$BASE_DN" -w "$ADMIN_PASS" -f /tmp/add-thomas-to-admin-keycloak.ldif || echo "[init_ldap] Erreur lors de l'ajout de thomas au groupe admin_keycloak"
echo "[init_ldap] Thomas ajouté au groupe admin_keycloak"

# Tentative de suppression d'un éventuel entry cn=admin (peu probable qu'il existe)
if ldapsearch -x -H ldap:/// -D "cn=admin,$BASE_DN" -w "$ADMIN_PASS" -b "$BASE_DN" "(cn=admin)" dn | grep -q "^dn: cn=admin,$BASE_DN$"; then
  echo "[init_ldap] Suppression cn=admin pour discrétisation..."
  cat > /tmp/delete-admin.ldif <<EOF
version: 1

dn: cn=admin,$BASE_DN
changetype: delete
EOF
  ldapmodify -x -H ldap:/// -D "cn=admin,$BASE_DN" -w "$ADMIN_PASS" -f /tmp/delete-admin.ldif || echo "[init_ldap] Erreur lors de la suppression de cn=admin"
fi

fi # fin rôle provider (DIT)

# Configuration des ACL de base pour permettre l'accès
echo "[init_ldap] Configuration des ACL de base..."
cat > /tmp/acl-basic.ldif <<EOF
dn: olcDatabase={1}mdb,cn=config
changetype: modify
add: olcAccess
olcAccess: to * by dn.exact=cn=admin,$BASE_DN manage by * read
EOF
ldapmodify -H ldapi:/// -Y EXTERNAL -f /tmp/acl-basic.ldif || echo "[init_ldap] Erreur lors de la configuration des ACL de base"

# Configuration des ACL avancées pour le groupe admin_ldap
echo "[init_ldap] Configuration des ACL avancées pour admin_ldap..."
cat > /tmp/acl-advanced.ldif <<EOF
dn: olcDatabase={1}mdb,cn=config
changetype: modify
add: olcAccess
olcAccess: to attrs=userPassword by group.exact=cn=admin_ldap,ou=groups,$BASE_DN write by self write by anonymous auth by * none
olcAccess: to dn.base="" by * read
olcAccess: to * by group.exact=cn=admin_ldap,ou=groups,$BASE_DN manage by self write by users read by * none
EOF
ldapmodify -H ldapi:/// -Y EXTERNAL -f /tmp/acl-advanced.ldif || echo "[init_ldap] Erreur lors de la configuration des ACL avancées"

echo "[init_ldap] Structure DIT et ACL configurées."
\end{verbatim}
}

\section{Script \texttt{scripts/init\_ldap\_linux\_integration.sh}}\label{ann:init-linux}
{\footnotesize
\begin{verbatim}
#!/usr/bin/env bash
# Script d'initialisation de l'intégration LDAP-Linux
# Ce script configure PAM, SSSD et NSS pour l'authentification LDAP

# Variables d'environnement
BASE_DN=${LDAP_BASE_DN:-"dc=example,dc=org"}
DOMAIN=${LDAP_DOMAIN:-"example.org"}
ADMIN_PASS=${LDAP_ADMIN_PASSWORD:-"admin"}

if [ "${LDAP_SERVICE_ROLE:-provider}" = "consumer" ]; then
  echo "[init_ldap_linux] Ignoré sur réplica (LDAP_SERVICE_ROLE=consumer)."
  exit 0
fi

if [ "${LDAP_SERVICE_ROLE:-provider}" = "meta" ]; then
  echo "[init_ldap_linux] Ignoré sur méta-annuaire (LDAP_SERVICE_ROLE=meta)."
  exit 0
fi

echo "[init_ldap_linux] Début de la configuration de l'intégration LDAP-Linux..."

# Attendre que slapd soit prêt
for i in {1..30}; do
  if ldapwhoami -H ldapi:/// -Y EXTERNAL >/dev/null 2>&1; then
    break
  fi
  sleep 1
done

# Configuration de NSS (Name Service Switch)
echo "[init_ldap_linux] Configuration de NSS..."
cat > /etc/nsswitch.conf <<EOF
# /etc/nsswitch.conf
# Configuration pour l'intégration LDAP

passwd:         files ldap
group:          files ldap
shadow:         files ldap
gshadow:        files

hosts:          files dns
networks:       files

protocols:      db files
services:       db files
ethers:         db files
rpc:            db files

netgroup:       ldap
EOF

# Configuration de libnss-ldap
echo "[init_ldap_linux] Configuration de libnss-ldap..."
cat > /etc/ldap/ldap.conf <<EOF
# Configuration LDAP pour NSS
BASE $BASE_DN
URI ldap://localhost:389
TLS_REQCERT never
EOF

# Configuration de PAM pour LDAP
echo "[init_ldap_linux] Configuration de PAM..."
cat > /etc/pam.d/common-auth <<EOF
# Configuration PAM pour l'authentification LDAP
auth    sufficient      pam_ldap.so
auth    sufficient      pam_unix.so nullok_secure use_first_pass
auth    required        pam_deny.so
EOF

cat > /etc/pam.d/common-account <<EOF
# Configuration PAM pour la gestion des comptes
account sufficient      pam_ldap.so
account sufficient      pam_unix.so
account required        pam_deny.so
EOF

cat > /etc/pam.d/common-password <<EOF
# Configuration PAM pour les mots de passe
password        sufficient      pam_ldap.so
password        sufficient      pam_unix.so nullok obscure min=4 max=8 md5
password        required        pam_deny.so
EOF

cat > /etc/pam.d/common-session <<EOF
# Configuration PAM pour les sessions
session required        pam_mkhomedir.so skel=/etc/skel umask=0022
session sufficient      pam_ldap.so
session sufficient      pam_unix.so
session required        pam_deny.so
EOF

# Configuration de nslcd (plus simple et fiable que SSSD dans un conteneur)
echo "[init_ldap_linux] Configuration de nslcd..."
cat > /etc/nslcd.conf <<EOF
# Configuration nslcd pour LDAP
uri ldap://localhost:389
base $BASE_DN
ldap_version 3
binddn cn=admin,$BASE_DN
bindpw $ADMIN_PASS
filter passwd (objectClass=posixAccount)
filter group (objectClass=posixGroup)
filter shadow (objectClass=shadowAccount)
map passwd homeDirectory homeDirectory
map passwd uidNumber uidNumber
map passwd gidNumber gidNumber
map passwd loginShell loginShell
map passwd gecos gecos
map group gidNumber gidNumber
map group memberUid memberUid
EOF

# Permissions pour nslcd
chmod 600 /etc/nslcd.conf
chown root:root /etc/nslcd.conf

# Configuration de libpam-ldap
echo "[init_ldap_linux] Configuration de libpam-ldap..."
cat > /etc/ldap.conf <<EOF
# Configuration pour libpam-ldap
base $BASE_DN
uri ldap://localhost:389
ldap_version 3
rootbinddn cn=admin,$BASE_DN
binddn cn=admin,$BASE_DN
bindpw $ADMIN_PASS
pam_password crypt
nss_base_passwd ou=people,$BASE_DN
nss_base_shadow ou=people,$BASE_DN
nss_base_group ou=groups,$BASE_DN
nss_map_objectclass posixAccount user
nss_map_objectclass shadowAccount user
nss_map_objectclass posixGroup group
nss_map_attribute uid sAMAccountName
nss_map_attribute homeDirectory unixHomeDirectory
nss_map_attribute shadowLastChange pwdLastSet
nss_map_objectclass posixAccount user
nss_map_objectclass shadowAccount user
nss_map_objectclass posixGroup group
nss_map_attribute uid sAMAccountName
nss_map_attribute homeDirectory unixHomeDirectory
nss_map_attribute shadowLastChange pwdLastSet
pam_login_attribute uid
pam_member_attribute member
pam_filter objectclass=inetOrgPerson
pam_password exop
EOF

# Ajout des attributs POSIX aux utilisateurs LDAP
echo "[init_ldap_linux] Ajout des attributs POSIX aux utilisateurs..."
cat > /tmp/add_posix_attrs.ldif <<EOF
dn: uid=thomas,ou=people,$BASE_DN
changetype: modify
add: objectClass
objectClass: posixAccount
objectClass: shadowAccount
-
add: uidNumber
uidNumber: 1001
-
add: gidNumber
gidNumber: 1001
-
add: homeDirectory
homeDirectory: /home/thomas
-
add: loginShell
loginShell: /bin/bash
-
add: gecos
gecos: Thomas Thomas

dn: uid=john,ou=people,$BASE_DN
changetype: modify
add: objectClass
objectClass: posixAccount
objectClass: shadowAccount
-
add: uidNumber
uidNumber: 1002
-
add: gidNumber
gidNumber: 1002
-
add: homeDirectory
homeDirectory: /home/john
-
add: loginShell
loginShell: /bin/bash
-
add: gecos
gecos: John John
EOF

ldapmodify -x -H ldap:/// -D "cn=admin,$BASE_DN" -w "$ADMIN_PASS" -f /tmp/add_posix_attrs.ldif || echo "[init_ldap_linux] Erreur lors de l'ajout des attributs POSIX"

# Conversion des groupes de groupOfNames à posixGroup
echo "[init_ldap_linux] Conversion des groupes en posixGroup..."
# Vérifier d'abord si les groupes existent et ont des membres
# Supprimer d'abord les membres existants (s'ils existent)
cat > /tmp/convert_groups_to_posix.ldif <<EOF
dn: cn=admin_ldap,ou=groups,$BASE_DN
changetype: modify
delete: objectClass
objectClass: groupOfNames
-
delete: member
EOF

# Ajouter les membres s'ils existent
if ldapsearch -x -H ldap:/// -D "cn=admin,$BASE_DN" -w "$ADMIN_PASS" -b "cn=admin_ldap,ou=groups,$BASE_DN" "(member=*)" member 2>/dev/null | grep -q "member:"; then
    # Supprimer tous les membres existants
    ldapsearch -x -H ldap:/// -D "cn=admin,$BASE_DN" -w "$ADMIN_PASS" -b "cn=admin_ldap,ou=groups,$BASE_DN" "(member=*)" member 2>/dev/null | grep "^member:" | sed 's/^member: /delete: member\nmember: /' >> /tmp/convert_groups_to_posix.ldif
fi

cat >> /tmp/convert_groups_to_posix.ldif <<EOF
-
add: objectClass
objectClass: posixGroup
-
add: gidNumber
gidNumber: 1001
-
add: memberUid
memberUid: thomas

dn: cn=developers,ou=groups,$BASE_DN
changetype: modify
delete: objectClass
objectClass: groupOfNames
-
delete: member
EOF

if ldapsearch -x -H ldap:/// -D "cn=admin,$BASE_DN" -w "$ADMIN_PASS" -b "cn=developers,ou=groups,$BASE_DN" "(member=*)" member 2>/dev/null | grep -q "member:"; then
    ldapsearch -x -H ldap:/// -D "cn=admin,$BASE_DN" -w "$ADMIN_PASS" -b "cn=developers,ou=groups,$BASE_DN" "(member=*)" member 2>/dev/null | grep "^member:" | sed 's/^member: /delete: member\nmember: /' >> /tmp/convert_groups_to_posix.ldif
fi

cat >> /tmp/convert_groups_to_posix.ldif <<EOF
-
add: objectClass
objectClass: posixGroup
-
add: gidNumber
gidNumber: 1002
-
add: memberUid
memberUid: john
EOF

ldapmodify -x -H ldap:/// -D "cn=admin,$BASE_DN" -w "$ADMIN_PASS" -f /tmp/convert_groups_to_posix.ldif 2>&1 || {
    echo "[init_ldap_linux] Tentative de conversion des groupes..."
    # Si la conversion échoue, essayer une approche plus simple : supprimer et recréer
    echo "[init_ldap_linux] Suppression et recréation des groupes en posixGroup..."
    # Supprimer les groupes existants
    ldapdelete -x -H ldap:/// -D "cn=admin,$BASE_DN" -w "$ADMIN_PASS" "cn=admin_ldap,ou=groups,$BASE_DN" 2>/dev/null || true
    ldapdelete -x -H ldap:/// -D "cn=admin,$BASE_DN" -w "$ADMIN_PASS" "cn=developers,ou=groups,$BASE_DN" 2>/dev/null || true
    sleep 1
    # Recréer les groupes en posixGroup
    cat > /tmp/recreate_posix_groups.ldif <<EOF
dn: cn=admin_ldap,ou=groups,$BASE_DN
objectClass: top
objectClass: posixGroup
cn: admin_ldap
gidNumber: 1001
memberUid: thomas

dn: cn=developers,ou=groups,$BASE_DN
objectClass: top
objectClass: posixGroup
cn: developers
gidNumber: 1002
memberUid: john
EOF
    ldapadd -x -H ldap:/// -D "cn=admin,$BASE_DN" -w "$ADMIN_PASS" -f /tmp/recreate_posix_groups.ldif 2>&1 && echo "[init_ldap_linux] Groupes recréés en posixGroup" || echo "[init_ldap_linux] Erreur lors de la recréation des groupes"
}

# Démarrage des services
echo "[init_ldap_linux] Démarrage des services..."
# Démarrer nslcd
service nslcd restart || service nslcd start
# Attendre que nslcd soit prêt
sleep 2
# Redémarrer nscd pour recharger la configuration et vider le cache
service nscd restart || service nscd start
# Vider le cache nscd pour forcer la relecture depuis LDAP
nscd -i passwd 2>/dev/null || true
nscd -i group 2>/dev/null || true
sleep 1

# Test de l'intégration
echo "[init_ldap_linux] Test de l'intégration..."
echo "Test de résolution des utilisateurs:"
getent passwd thomas
getent passwd john

echo "Test de résolution des groupes:"
getent group admin_ldap
getent group developers

echo "[init_ldap_linux] Configuration de l'intégration LDAP-Linux terminée."
\end{verbatim}
}

\section{Script \texttt{scripts/init\_replication\_provider.sh}}\label{ann:repl-provider}
{\footnotesize
\begin{verbatim}
#!/usr/bin/env bash
# Fournisseur de réplication : overlay syncprov + olcServerID (OpenLDAP 2.6).
set -uo pipefail

if [ "${LDAP_SERVICE_ROLE:-provider}" = "consumer" ]; then
  exit 0
fi

if [ "${LDAP_SERVICE_ROLE:-provider}" = "meta" ]; then
  exit 0
fi

for _ in $(seq 1 30); do
  if ldapwhoami -H ldapi:/// -Y EXTERNAL >/dev/null 2>&1; then
    break
  fi
  sleep 1
done

if ldapsearch -H ldapi:/// -Y EXTERNAL -b cn=config -s sub -LLL '(objectClass=olcSyncProvConfig)' dn 2>/dev/null | grep -q '^dn:'; then
  echo "[init_replication_provider] Overlay syncprov déjà présent."
else
  echo "[init_replication_provider] Chargement module syncprov.la sur cn=module{0}..."
  ldapmodify -H ldapi:/// -Y EXTERNAL <<'EOF' 2>/dev/null || true
dn: cn=module{0},cn=config
changetype: modify
add: olcModuleLoad
olcModuleLoad: syncprov.la
EOF
  echo "[init_replication_provider] Ajout overlay syncprov sur mdb {1}..."
  ldapadd -H ldapi:/// -Y EXTERNAL <<'EOF'
dn: olcOverlay=syncprov,olcDatabase={1}mdb,cn=config
objectClass: olcOverlayConfig
objectClass: olcSyncProvConfig
olcOverlay: syncprov
EOF
fi

if ldapsearch -H ldapi:/// -Y EXTERNAL -b cn=config -s base -LLL olcServerID 2>/dev/null | grep -q "^olcServerID: 1"; then
  echo "[init_replication_provider] olcServerID=1 déjà présent."
else
  echo "[init_replication_provider] Ajout olcServerID=1..."
  ldapmodify -H ldapi:/// -Y EXTERNAL <<'EOF' 2>/dev/null || true
dn: cn=config
changetype: modify
add: olcServerID
olcServerID: 1
EOF
fi

echo "[init_replication_provider] Index entryUUID (recommandé pour syncrepl)..."
ldapmodify -H ldapi:/// -Y EXTERNAL <<'EOF' 2>/dev/null || true
dn: olcDatabase={1}mdb,cn=config
changetype: modify
add: olcDbIndex
olcDbIndex: entryUUID eq
EOF

echo "[init_replication_provider] Terminé."
\end{verbatim}
}

\section{Script \texttt{scripts/init\_replication\_consumer.sh}}\label{ann:repl-consumer}
{\footnotesize
\begin{verbatim}
#!/usr/bin/env bash
# Réplica en lecture seule : olcSyncrepl + olcReadOnly (les écritures client sont refusées, pas la synchro).
set -uo pipefail

if [ "${LDAP_SERVICE_ROLE:-provider}" != "consumer" ]; then
  exit 0
fi

BASE_DN=${LDAP_BASE_DN:-dc=example,dc=org}
ADMIN_PASS=${LDAP_ADMIN_PASSWORD:-admin}
PROVIDER_URI=${LDAP_REPLICATION_PROVIDER_URI:-ldap://ldap:389}

for _ in $(seq 1 30); do
  if ldapwhoami -H ldapi:/// -Y EXTERNAL >/dev/null 2>&1; then
    break
  fi
  sleep 1
done

echo "[init_replication_consumer] Attente du fournisseur ${PROVIDER_URI}..."
ready=0
for _ in $(seq 1 90); do
  if ldapsearch -x -H "$PROVIDER_URI" -D "cn=admin,$BASE_DN" -w "$ADMIN_PASS" -b "$BASE_DN" -s base dn >/dev/null 2>&1 \
    && ldapsearch -x -H "$PROVIDER_URI" -D "cn=admin,$BASE_DN" -w "$ADMIN_PASS" -b "ou=people,$BASE_DN" -s base dn >/dev/null 2>&1; then
    ready=1
    break
  fi
  sleep 2
done
if [ "$ready" != 1 ]; then
  echo "[init_replication_consumer] AVERTISSEMENT : fournisseur peu joignable ; syncrepl pourrait échouer au démarrage."
fi

if ! ldapsearch -H ldapi:/// -Y EXTERNAL -b cn=config -s base -LLL olcServerID 2>/dev/null | grep -q "^olcServerID: 2"; then
  echo "[init_replication_consumer] olcServerID=2 sur cn=config..."
  ldapmodify -H ldapi:/// -Y EXTERNAL <<'EOF' 2>/dev/null || true
dn: cn=config
changetype: modify
add: olcServerID
olcServerID: 2
EOF
fi

if ldapsearch -H ldapi:/// -Y EXTERNAL -b olcDatabase={1}mdb,cn=config -s base -LLL olcSyncrepl 2>/dev/null | grep -q "^olcSyncrepl:"; then
  echo "[init_replication_consumer] olcSyncrepl déjà configuré."
  exit 0
fi

echo "[init_replication_consumer] Configuration olcSyncrepl + olcReadOnly..."
ldapmodify -H ldapi:/// -Y EXTERNAL <<EOF
dn: olcDatabase={1}mdb,cn=config
changetype: modify
add: olcReadOnly
olcReadOnly: TRUE
-
add: olcSyncrepl
olcSyncrepl: rid=001 provider=${PROVIDER_URI} bindmethod=simple binddn="cn=admin,${BASE_DN}" credentials=${ADMIN_PASS} searchbase="${BASE_DN}" type=refreshAndPersist retry="60 +" timeout=1 schemachecking=off scope=sub
EOF

echo "[init_replication_consumer] Terminé."
\end{verbatim}
}

\section{Script \texttt{scripts/init\_meta\_annuaire.sh}}\label{ann:init-meta}
\noindent\small\textit{Dans l'image Docker : copié en \texttt{/container/init\_meta\_annuaire.sh} (hors \texttt{init.d}), appelé par \texttt{init\_ldap.sh} si \texttt{LDAP\_SERVICE\_ROLE=meta}.}\par\medskip
{\footnotesize
\begin{verbatim}
#!/usr/bin/env bash
# Méta-annuaire OpenLDAP (back-meta) : agrège plusieurs annuaires via suffixmassage.
# Invoqué depuis init_ldap.sh lorsque LDAP_SERVICE_ROLE=meta.
set -euo pipefail

ADMIN_PASS=${LDAP_ADMIN_PASSWORD:-admin}
META_SUFFIX=${LDAP_META_SUFFIX:-o=federation}
META_ROOT_DN="cn=admin,${META_SUFFIX}"

EX_URI=${LDAP_META_EXAMPLE_URI:-ldap://ldap:389}
EX_REAL=${LDAP_META_EXAMPLE_SUFFIX:-dc=example,dc=org}
EX_VIRT=${LDAP_META_EXAMPLE_VIRTUAL:-ou=example-org,o=federation}

AC_URI=${LDAP_META_ACME_URI:-ldap://ldap-acme:389}
AC_REAL=${LDAP_META_ACME_SUFFIX:-dc=acme,dc=com}
AC_VIRT=${LDAP_META_ACME_VIRTUAL:-ou=acme-corp,o=federation}

export LDAPTLS_REQCERT=never

for _ in $(seq 1 30); do
  if ldapwhoami -H ldapi:/// -Y EXTERNAL >/dev/null 2>&1; then
    break
  fi
  sleep 1
done

echo "[init_meta] Suppression de la base mdb par défaut..."
ldapmodify -H ldapi:/// -Y EXTERNAL <<'EOF' || true
dn: olcDatabase={1}mdb,cn=config
changetype: delete
EOF

echo "[init_meta] Purge des fichiers mdb et redémarrage de slapd..."
pkill -u openldap -TERM slapd 2>/dev/null || pkill -TERM slapd 2>/dev/null || true
for _ in $(seq 1 40); do
  pgrep slapd >/dev/null 2>&1 || break
  sleep 0.25
done
pkill -u openldap -KILL slapd 2>/dev/null || pkill -KILL slapd 2>/dev/null || true
find /var/lib/ldap -mindepth 1 -delete 2>/dev/null || true
chown openldap:openldap /var/lib/ldap
slapd -h "ldap:/// ldapi:///" -u openldap -g openldap -F /etc/ldap/slapd.d &
for _ in $(seq 1 60); do
  if ldapwhoami -H ldapi:/// -Y EXTERNAL >/dev/null 2>&1; then
    break
  fi
  sleep 1
done

if ! ldapsearch -H ldapi:/// -Y EXTERNAL -b cn=module{0},cn=config -s base -LLL olcModuleLoad 2>/dev/null | grep -qF back_ldap; then
  echo "[init_meta] Chargement des modules back_ldap (requis par back_meta) puis back_meta..."
  ldapmodify -H ldapi:/// -Y EXTERNAL <<'EOF' || true
dn: cn=module{0},cn=config
changetype: modify
add: olcModuleLoad
olcModuleLoad: back_ldap.la
-
add: olcModuleLoad
olcModuleLoad: back_meta.la
EOF
fi

HASH=$(slappasswd -s "$ADMIN_PASS")
echo "[init_meta] Création de la base meta (suffix=$META_SUFFIX)..."
ldapadd -H ldapi:/// -Y EXTERNAL <<EOF
dn: olcDatabase=meta,cn=config
objectClass: olcDatabaseConfig
objectClass: olcMetaConfig
olcDatabase: meta
olcSuffix: $META_SUFFIX
olcRootDN: $META_ROOT_DN
olcRootPW: $HASH
olcDbOnErr: continue
olcDbCancel: abandon
olcDbTFSupport: no
olcAccess: to * by dn.exact=$META_ROOT_DN manage by * read
EOF

META_DB_DN=$(ldapsearch -H ldapi:/// -Y EXTERNAL -b cn=config -s one -LLL '(olcDatabase=meta)' dn 2>/dev/null | sed -n 's/^dn: //p' | head -1)
if [ -z "$META_DB_DN" ]; then
  echo "[init_meta] ERREUR : entrée olcDatabase meta introuvable dans cn=config."
  exit 1
fi

# Le namingContext dans olcDbURI doit être un sous-arbre du suffixe meta (olcSuffix), pas le suffixe réel du backend.
echo "[init_meta] Cible meta : $EX_VIRT → $EX_URI + suffixmassage → $EX_REAL"
ldapadd -H ldapi:/// -Y EXTERNAL <<EOF
dn: olcMetaSub={0}example,$META_DB_DN
objectClass: olcConfig
objectClass: olcMetaTargetConfig
olcMetaSub: {0}example
olcDbURI: $EX_URI/$EX_VIRT
olcDbRewrite: {0}suffixmassage "$EX_VIRT" "$EX_REAL"
olcDbIDAssertBind: bindmethod=simple binddn="cn=admin,$EX_REAL" credentials=$ADMIN_PASS
olcDbChaseReferrals: FALSE
EOF

echo "[init_meta] Cible meta : $AC_VIRT → $AC_URI + suffixmassage → $AC_REAL"
ldapadd -H ldapi:/// -Y EXTERNAL <<EOF
dn: olcMetaSub={1}acme,$META_DB_DN
objectClass: olcConfig
objectClass: olcMetaTargetConfig
olcMetaSub: {1}acme
olcDbURI: $AC_URI/$AC_VIRT
olcDbRewrite: {0}suffixmassage "$AC_VIRT" "$AC_REAL"
olcDbIDAssertBind: bindmethod=simple binddn="cn=admin,$AC_REAL" credentials=$ADMIN_PASS
olcDbChaseReferrals: FALSE
EOF

echo "[init_meta] Méta-annuaire configuré ($META_SUFFIX)."
exit 0
\end{verbatim}
}

\section{Script \texttt{scripts/configure\_keycloak\_ldap.sh} (machine hôte)}\label{ann:keycloak-sh}
\noindent\small\textit{Hors image : à exécuter sur la machine hôte après \texttt{docker compose up -d}. Fichier dans le dépôt : \texttt{projet/scripts/configure\_keycloak\_ldap.sh} (ex.\ \texttt{bash projet/scripts/configure\_keycloak\_ldap.sh} depuis la racine du clone, ou \texttt{bash scripts/configure\_keycloak\_ldap.sh} depuis \texttt{projet/}). Pas de second listing dans le corps du guide.}\par\medskip
{\footnotesize
\begin{verbatim}
#!/usr/bin/env bash
# Configure Keycloak : realm + User Federation LDAP (API d'administration).
# À lancer depuis la machine hôte une fois « docker compose up -d » (Keycloak joindra ldap:389).
set -euo pipefail

KEYCLOAK_URL="${KEYCLOAK_URL:-http://localhost:8090}"
KEYCLOAK_ADMIN="${KEYCLOAK_ADMIN:-admin}"
KEYCLOAK_ADMIN_PASSWORD="${KEYCLOAK_ADMIN_PASSWORD:-admin}"
REALM="${KEYCLOAK_REALM:-tp-ldap}"
LDAP_CONNECTION_URL="${KEYCLOAK_LDAP_CONNECTION:-ldap://ldap:389}"
LDAP_BIND_DN="${LDAP_BIND_DN:-cn=admin,dc=example,dc=org}"
LDAP_BIND_PW="${LDAP_BIND_PW:-admin}"
LDAP_USERS_DN="${LDAP_USERS_DN:-ou=people,dc=example,dc=org}"

echo "[configure_keycloak_ldap] Attente de Keycloak (${KEYCLOAK_URL})..."
ready=0
for _ in $(seq 1 90); do
  if curl -sf "${KEYCLOAK_URL}/realms/master" >/dev/null; then
    ready=1
    break
  fi
  sleep 2
done
if [ "$ready" != 1 ]; then
  echo "[configure_keycloak_ldap] ERREUR : Keycloak ne répond pas à temps."
  exit 1
fi

echo "[configure_keycloak_ldap] Obtention du jeton admin-cli..."
TOKEN_JSON=$(curl -sS -X POST "${KEYCLOAK_URL}/realms/master/protocol/openid-connect/token" \
  -H "Content-Type: application/x-www-form-urlencoded" \
  --data-urlencode "grant_type=password" \
  --data-urlencode "client_id=admin-cli" \
  --data-urlencode "username=${KEYCLOAK_ADMIN}" \
  --data-urlencode "password=${KEYCLOAK_ADMIN_PASSWORD}")

TOKEN=$(printf '%s' "$TOKEN_JSON" | python3 -c "import sys, json; print(json.load(sys.stdin).get('access_token',''))" 2>/dev/null || true)
if [ -z "$TOKEN" ]; then
  echo "[configure_keycloak_ldap] ERREUR : impossible d'obtenir access_token. Réponse :"
  printf '%s\n' "$TOKEN_JSON"
  exit 1
fi

AUTH_HDR=(curl -sS -H "Authorization: Bearer ${TOKEN}")
REALM_CODE=$("${AUTH_HDR[@]}" -o /dev/null -w "%{http_code}" "${KEYCLOAK_URL}/admin/realms/${REALM}")

if [ "$REALM_CODE" != "200" ]; then
  echo "[configure_keycloak_ldap] Création du realm « ${REALM} » (HTTP realm=${REALM_CODE})..."
  HTTP_CREATE=$("${AUTH_HDR[@]}" -o /dev/null -w "%{http_code}" -X POST "${KEYCLOAK_URL}/admin/realms" \
    -H "Content-Type: application/json" \
    -d "$(python3 -c "import json; print(json.dumps({'realm':'${REALM}','enabled':True,'displayName':'TP LDAP','registrationAllowed':False,'loginWithEmailAllowed':True,'duplicateEmailsAllowed':False}))")")
  if [ "$HTTP_CREATE" != "201" ]; then
    echo "[configure_keycloak_ldap] ERREUR : création realm HTTP ${HTTP_CREATE}"
    exit 1
  fi
else
  echo "[configure_keycloak_ldap] Realm « ${REALM} » déjà présent."
fi

EXISTING_ID=$("${AUTH_HDR[@]}" "${KEYCLOAK_URL}/admin/realms/${REALM}/components?type=org.keycloak.storage.UserStorageProvider" | python3 -c "
import sys, json
try:
    for c in json.load(sys.stdin):
        if c.get('providerId') == 'ldap':
            print(c.get('id', '') or '')
            break
except json.JSONDecodeError:
    pass
")

if [ -n "$EXISTING_ID" ]; then
  echo "[configure_keycloak_ldap] Fournisseur LDAP déjà configuré (id=${EXISTING_ID})."
  STORAGE_ID="$EXISTING_ID"
else
  echo "[configure_keycloak_ldap] Création du composant User Federation LDAP..."
  REALM_INTERNAL_ID=$("${AUTH_HDR[@]}" "${KEYCLOAK_URL}/admin/realms/${REALM}" | python3 -c "import sys, json; print(json.load(sys.stdin)['id'])")

  export _KC_REALM_INTERNAL_ID="$REALM_INTERNAL_ID"
  export _KC_LDAP_URL="$LDAP_CONNECTION_URL"
  export _KC_LDAP_BIND_DN="$LDAP_BIND_DN"
  export _KC_LDAP_BIND_PW="$LDAP_BIND_PW"
  export _KC_LDAP_USERS_DN="$LDAP_USERS_DN"
  BODY=$(python3 <<'PY'
import json, os
print(json.dumps({
  "name": "ldap",
  "providerId": "ldap",
  "providerType": "org.keycloak.storage.UserStorageProvider",
  "parentId": os.environ["_KC_REALM_INTERNAL_ID"],
  "config": {
    "enabled": ["true"],
    "priority": ["0"],
    "importEnabled": ["true"],
    "editMode": ["READ_ONLY"],
    "syncRegistrations": ["false"],
    "vendor": ["other"],
    "usernameLDAPAttribute": ["uid"],
    "rdnLDAPAttribute": ["uid"],
    "uuidLDAPAttribute": ["entryUUID"],
    "userObjectClasses": ["inetOrgPerson, organizationalPerson"],
    "connectionUrl": [os.environ["_KC_LDAP_URL"]],
    "bindDn": [os.environ["_KC_LDAP_BIND_DN"]],
    "bindCredential": [os.environ["_KC_LDAP_BIND_PW"]],
    "usersDn": [os.environ["_KC_LDAP_USERS_DN"]],
    "authType": ["simple"],
    "searchScope": ["2"],
    "pagination": ["true"],
    "connectionPooling": ["true"],
    "startTls": ["false"],
  },
}))
PY
)
  unset _KC_REALM_INTERNAL_ID _KC_LDAP_URL _KC_LDAP_BIND_DN _KC_LDAP_BIND_PW _KC_LDAP_USERS_DN
  CREATE_OUT=$(mktemp)
  HTTP_CODE=$(curl -sS -o "$CREATE_OUT" -w "%{http_code}" -X POST \
    "${KEYCLOAK_URL}/admin/realms/${REALM}/components" \
    -H "Authorization: Bearer ${TOKEN}" \
    -H "Content-Type: application/json" \
    -d "$BODY")
  if [ "$HTTP_CODE" != "201" ]; then
    echo "[configure_keycloak_ldap] ERREUR création LDAP provider (HTTP ${HTTP_CODE}) :"
    cat "$CREATE_OUT"
    rm -f "$CREATE_OUT"
    exit 1
  fi
  rm -f "$CREATE_OUT"
  STORAGE_ID=$("${AUTH_HDR[@]}" "${KEYCLOAK_URL}/admin/realms/${REALM}/components?type=org.keycloak.storage.UserStorageProvider" | python3 -c "
import sys, json
for c in json.load(sys.stdin):
    if c.get('providerId') == 'ldap':
        print(c['id'])
        break
")
  if [ -z "$STORAGE_ID" ]; then
    echo "[configure_keycloak_ldap] ERREUR : id du stockage LDAP introuvable après création."
    exit 1
  fi
fi

echo "[configure_keycloak_ldap] Synchronisation complète des utilisateurs LDAP..."
SYNC_CODE=$(curl -sS -o /dev/null -w "%{http_code}" -X POST \
  "${KEYCLOAK_URL}/admin/realms/${REALM}/user-storage/${STORAGE_ID}/sync?action=triggerFullSync" \
  -H "Authorization: Bearer ${TOKEN}")
if [ "$SYNC_CODE" != "204" ] && [ "$SYNC_CODE" != "200" ]; then
  echo "[configure_keycloak_ldap] AVERTISSEMENT : sync HTTP ${SYNC_CODE} (souvent acceptable si déjà synchronisé)."
fi

echo "[configure_keycloak_ldap] Terminé. Console : ${KEYCLOAK_URL}/admin/ - realm « ${REALM} »."
\end{verbatim}
}

